\chapter[Zusammenfassung (Summary in German)]{Zusammenfassung {\hfill\small (Summary in German)}}

\mnote{Set language settings in main.tex. I chose ``british'' in that file but ngerman here (date format: \today). If cleverref is working, the following word should be ``Kapitel II'': } \cref{ch:1_chapter}.

\mnote{Random text: } In jenem Versuche, das bisherige Verfahren der Metaphysik umzu\"andern,
und dadurch, dass wir nach dem Beispiele der Geometer und Naturforscher
eine g\"anzliche Revolution mit derselben vornehmen, besteht nun das
Gesch\"aft dieses Kritik der reinen spekulativen Vernunft. Sie ist ein
Traktat von der Methode, nicht ein System der Wissenschaft selbst;
aber sie verzeichnet gleichwohl den ganzen Umriss derselben, sowohl
in Ansehung ihrer Grenzen, als auch den ganzen inneren Gliederbau
derselben. 
%
	\begin{enumerate}[itemsep=0ex, topsep=1ex]% could also be set globally
	\item \mnote{See the code for the spacing of enumerates.}
	\item PS: Special symbols in the code should be ok (UTF8 encoding): ä ü ö ß @
	\item Check the siunitx package for the correct typesetting of dimensional quantities such as ``\SI{1}{\micro\litre}''.
\end{enumerate}
%
\mnote{Random text: } Denn das hat die reine, spekulative Vernunft Eigent\"umliches
an sich, dass sie ihr eigen Verm\"ogen, nach Verschiedenheit der Art,
wie sie sich Objekte zum Denken w\"ahlt, ausmessen, und auch selbst die
mancherlei Arten, sich Aufgaben vorzulegen, vollst\"andig vorz\"ahlen, und
so den ganzen Vorri\ss{} zu einem System der Metaphysik verzeichnen kann
und soll; weil, was das erste betrifft, in der Erkenntnis a priori den
Objekten nichts beigelegt werden kann, als was das denkende Subjekt
aus sich selbst hernimmt, und, was das zweite anlangt, sie in Ansehung
der Erkenntnisprinzipien eine ganz abgesonderte, f\"ur sich bestehende
Einheit ist, in welcher ein jedes Glied, wie in einem organisierten
K\"orper, um aller anderen und alle um eines willen da sind, und kein
Prinzip mit Sicherheit in einer Beziehung genommen werden kann,
ohne es zugleich in der durchg\"angigen Beziehung zum ganzen
reinen Vernunftgebrauch untersucht zu haben. 

Daf\"ur aber hat
auch die Metaphysik das seltene Gl\"uck, welches keiner anderen
Vernunftwissenschaft, die es mit Objekten zu tun hat (denn die Logik
besch\"aftigt sich nur mit der Form des Denkens \"uberhaupt), zuteil
werden kann, dass, wenn sie durch diese Kritik in den sicheren Gang
einer Wissenschaft gebracht worden, sie das ganze Feld der f\"ur sie
geh\"origen Erkenntnisse v\"ollig befassen und also ihr Werk vollenden
und f\"ur die Nachwelt, als einen nie zu vermehrenden Hauptstuhl.



% Just a gimmick: I put a fleuron on otherwise empty pages
\clearpage
\thispagestyle{empty}
\vspace*{30mm}
\begin{center}
\includegraphics{fleuron}
\end{center}
\clearpage


